\documentclass[11pt]{article}
\usepackage[margin=1.1in]{geometry}
\usepackage{times}
\usepackage[T1]{fontenc}
\usepackage{setspace}
\onehalfspacing
\usepackage{parskip}
\pagenumbering{gobble}

\title{Benchmark Leaderboards}
\author{Flyxion}
\date{September 2026}

\begin{document}
\maketitle

A laboratory builds a benchmark, trains near the tasks it measures, tunes against the surrounding evaluation culture, selects its strongest checkpoint, and announces the checkpoint has surpassed humanity.

Humanity, in this comparison, is represented by a handful of contractors answering unfamiliar questions under time pressure for an undisclosed rate. The model is represented by the sole survivor of an optimization process that cost several hundred million dollars and discarded every worse attempt along the way. Fairness, in this setup, requires only that neither side be permitted to complain, which is easy to arrange when one side is a spreadsheet.

A real exam works because the party who writes the questions, the party who takes them, and the party who grades them are three different parties, none able to adjust the others' incentives. Collapse those three into one and the leaderboard stops measuring the capability it was built to proxy for and starts measuring how well the model has been optimized against the proxy itself, which is a different quantity that happens to share a name with the first one. Contamination compounds it: a model that has seen the questions, or their near cousins, during training is not generalizing, it's recalling, in a format designed to look like the former. When the benchmark eventually saturates, this is reported as evidence that intelligence has been solved rather than evidence that the test has been memorized, and a new benchmark is introduced so that intelligence can be solved again next quarter, by the same method, against the same absence of anyone in the room whose job is to say the number went down.

\end{document}
